\documentclass[11pt]{article}

% ---------------------------------------------------------------
%  AI Applications in Research: The Use of Claude Code
%  A conceptual summary synthesised from Scott Cunningham's
%  "Scott's Mixtape" Claude Code essay series.
% ---------------------------------------------------------------

\usepackage[utf8]{inputenc}
\usepackage[T1]{fontenc}
\usepackage[margin=1in]{geometry}
\usepackage{enumitem}
\usepackage{titlesec}
\usepackage{hyperref}
\usepackage{xcolor}

\hypersetup{
    colorlinks=true,
    linkcolor=black,
    urlcolor=blue!60!black,
    citecolor=black
}

\titleformat{\section}{\large\bfseries}{\thesection}{1em}{}
\titleformat{\subsection}{\normalsize\bfseries}{\thesubsection}{1em}{}

\setlist[itemize]{itemsep=2pt, topsep=3pt}

\title{\textbf{AI Applications in Research}\\[2pt]
\large Core Concepts and the Use of Claude Code in Empirical Work}
\author{Summary compiled from Scott Cunningham's \emph{Scott's Mixtape} essay series}
\date{July 2026}

\begin{document}

\maketitle

\begin{abstract}
\noindent This document summarises the main concepts surrounding the use of
AI agents---and specifically \emph{Claude Code}---in academic and empirical
research. It draws on the themes developed across Scott Cunningham's
long-running essay series on the subject, which documents in public the
practical, economic, and pedagogical consequences of integrating agentic AI
into the daily work of an applied econometrician. The goal is to give a
scannable, conceptual map: what these tools change, where the new bottlenecks
lie, and how researchers can adopt them responsibly.
\end{abstract}

\section{Introduction}

Agentic coding assistants such as Claude Code have moved AI from a
conversational novelty to a genuine part of the research \emph{workflow}. Rather
than answering isolated questions, an agent can read a codebase, run analyses,
inspect data, replicate published results, and iterate toward a verified answer.
For empirical researchers this represents a shift in how work is produced, not
merely a faster way of typing. The sections below organise the recurring
concepts into five clusters: the economics of research production, practical
workflows, verification and quality control, pedagogy, and the risks that
accompany adoption.

\section{The Economics of Research Production}

The most consequential ideas are economic. When the marginal cost of producing
analysis falls sharply, both the volume of output and the profile of errors
change.

\begin{itemize}
  \item \textbf{Falling marginal costs of production.} Once an agent can carry
        out a large share of the mechanical labour of empirical work, the cost
        of producing an additional analysis, table, or replication drops toward
        zero, dramatically increasing potential output.
  \item \textbf{Aggregate errors scale with output.} Cheaper production means
        \emph{more} production, and if per-task error rates stay constant, the
        total number of errors in circulation can rise even as individual work
        improves---making quality control a first-order concern.
  \item \textbf{The zero-profit condition.} As the tools diffuse and the cost of
        competent empirical output collapses, the ``economic rent'' that once
        accrued to technical skill alone is competed away, pushing the source of
        value elsewhere (judgement, questions, verification, credibility).
  \item \textbf{Research and publishing are diverging.} The act of \emph{doing}
        research and the act of \emph{publishing} it are becoming distinct
        activities, with different bottlenecks, timelines, and value.
\end{itemize}

\section{Verification as the New Bottleneck}

A central and repeated theme is that generation is no longer the constraint---
\emph{trust} is.

\begin{itemize}
  \item \textbf{Verification is the new bottleneck.} When output is cheap and
        abundant, the scarce input becomes confidence that the output is
        correct; the researcher's time shifts from producing to checking.
  \item \textbf{Checklists for agentic research.} Reliable applied econometrics
        with AI agents benefits from explicit, staged checklists (design,
        estimation, robustness) that make each analytical decision auditable
        rather than implicit.
  \item \textbf{Multi-agent audits.} One agent's work can be independently
        checked by another, adversarial pass---replicating and cross-examining
        results (e.g.\ auditing an estimator implementation) to surface errors a
        single run would miss.
  \item \textbf{Replication as method.} Reproducing published results end-to-end
        becomes a practical, teachable exercise---a way both to validate the
        tooling and to learn the underlying method deeply.
\end{itemize}

\section{Practical Workflows and Tooling}

Alongside the ``big think'' economics runs a stream of concrete, reusable
practice---the ``how-to'' layer that makes adoption real.

\begin{itemize}
  \item \textbf{Skills and reusable prompts.} Packaging repeated tasks (slide
        decks, PDF handling, formatting) as named, reusable skills turns one-off
        prompting into durable infrastructure.
  \item \textbf{MCP integrations.} The Model Context Protocol lets an agent
        connect to external services and data sources---so the tool can act on a
        researcher's actual corpus, not just answer in the abstract.
  \item \textbf{Harness principles.} Thinking carefully about the ``harness''---
        how the agent is given context, tools, and guardrails---matters as much
        as the model itself for getting dependable results.
  \item \textbf{Artifacts as deliverables.} Increasingly the output itself
        (a deck, a video, a formatted document) \emph{is} the product, produced
        directly by the agent.
\end{itemize}

\section{Pedagogy and Integration into Teaching}

A distinctive concept is that these tools do not just change research---they
change how research \emph{skills} are taught.

\begin{itemize}
  \item \textbf{Teaching AI-agentic research.} Graduate curricula are beginning
        to pair deep methodological knowledge (econometrics, causal inference)
        with applied competence in \emph{using} AI agents to do research.
  \item \textbf{A textbook drafted in public.} Writing openly and iteratively
        about the tools---``thinking out loud''---doubles as pedagogy and as raw
        material for structured teaching resources.
  \item \textbf{Documenting progress transparently.} Keeping an ongoing,
        honest record of experiments, failures, and audits builds a body of
        material others can learn from and reproduce.
\end{itemize}

\section{Risks and Open Questions}

Adoption is presented as a genuine double-edged sword, not an unambiguous win.

\begin{itemize}
  \item \textbf{Amnesia and lost context.} Agents do not persist memory across
        sessions by default, so continuity, context, and reproducibility must be
        engineered deliberately.
  \item \textbf{Volume without discipline.} The ease of producing analysis can
        encourage over-production and, without checklists and verification,
        undermine rigour.
  \item \textbf{Erosion of technical rents.} As competent output becomes cheap,
        researchers must find new sources of distinctive value---questions,
        judgement, credibility---rather than technical execution alone.
  \item \textbf{The essayist's discipline.} Careful writing and reasoning remain
        the mechanism by which a researcher actually thinks a problem through;
        the tool accelerates execution but does not replace understanding.
\end{itemize}

\section{Conclusion}

The recurring message across the series is pragmatic optimism tempered by
caution. Claude Code and similar agents collapse the marginal cost of empirical
production and, in doing so, relocate the binding constraint from \emph{doing}
the work to \emph{trusting} it. The researcher's comparative advantage moves
upward: toward asking the right questions, designing auditable workflows,
verifying results, and teaching others to do the same. Handled with explicit
checklists, independent verification, and honest documentation, these tools
expand what a single researcher can credibly accomplish; handled carelessly,
they simply multiply errors faster than before.

\vspace{1em}
\noindent\rule{\linewidth}{0.4pt}\\[4pt]
\small\emph{Source material:} Scott Cunningham, \emph{Scott's Mixtape} Substack,
Claude Code essay series (Dec.\ 2025--Jul.\ 2026),
\href{https://causalinf.substack.com/}{causalinf.substack.com}.

\end{document}
